\documentclass[10pt,twocolumn,letterpaper]{article}
\usepackage[utf8]{inputenc}
\usepackage{geometry}
\geometry{letterpaper, margin=1in}
\usepackage{amsmath,amssymb,amsfonts}
\usepackage{booktabs}
\usepackage{microtype}
\usepackage{cite}
\usepackage{appendix}

\title{\textbf{Speculative KV Cache Quantization for Agentic Serving}}
\author{
    Arihan Iyer \\ \texttt{iamarihan@umich.edu} \and
    Kushagra Mishra \\ \texttt{kushm@umich.edu} \and
    Jatin Shihora \\ \texttt{jatinsh@umich.edu} \and
    Rohit Puthucode \\ \texttt{rohitps@umich.edu}
}

\date{}

\begin{document}

\maketitle

\begin{abstract}
The key-value (KV) cache is a foundational mechanism for accelerating autoregressive generation in Large Language Models (LLMs). However, its linear growth imposes severe GPU memory pressure in long-context and multi-turn agentic environments. Existing approximation techniques rely either on  cache eviction (which drops tokens permanently) or uniform low-bit quantization (which degrades critical contextual nuances). Recent draft-based approaches incorporate future lookahead to guide eviction, but they remain confined to single greedy generation trajectories. We propose a speculative multi-rollout KV cache quantization framework. When GPU memory pressure crosses a critical threshold, our system launches $N$ parallel lookahead rollouts using an inexpensive draft model to anticipate diverse future execution paths. By extracting consensus attention signals across rollouts, the system isolates high-value tokens to retain in high precision  while compressing the remaining context into a low-bit representation. This ensures that memory footprint is significantly reduced while preventing unrecoverable contextual degradation.
\end{abstract}

\section{Introduction}
Modern LLM agents frequently operate over long horizons, executing multi-turn workflows that incorporate extensive system instructions, tool documentation, retrieval chunks, and execution histories. In these settings, context sizes rapidly expand toward $100\text{k}$ to $1\text{M}$ tokens. Storing the KV cache for such sequences consumes massive High-Bandwidth Memory (HBM). For instance, serving a 32B-parameter model on an 80\,GB GPU leaves limited remaining memory for concurrent KV states, directly constraining batch sizes and serving throughput.

To alleviate memory pressure, prior work explores post-training KV cache reduction. Cache eviction policies (e.g., H2O\cite{h2023}, SnapKV\cite{snapkv24}) permanently discard tokens considered unimportant based on past attention scores. However, in agentic workloads with non-linear or shifting execution goals, discarded tokens often become critical in subsequent steps, leading to unrecoverable reasoning errors. Conversely, uniform post-training quantization (e.g., KIVI\cite{kivi24}) compresses all tokens uniformly, preserving sequence integrity but risking numerical degradation on outlier channels and sensitive attention sinks.

Recent work, such as SpecKV \cite{speckv26}, demonstrates that looking ahead into future generations via a small draft model yields significantly more accurate token-importance estimates than retrospective attention signals alone. Nonetheless, existing lookahead systems employ single greedy draft paths, making them vulnerable to draft hallucinations and branch divergence.

We propose an adaptive, speculative mixed-precision KV cache manager. Rather than dropping tokens permanently, our system dynamically quantizes tokens based on consensus attention importance derived from $N$ stochastic rollouts of a lightweight draft model. High-consensus tokens are pinned in full precision, while non-critical tokens are down-quantized to lower bit representations, achieving high compression without irreversible information loss.

\section{Related Work}
\noindent\textbf{KV Cache Compression vs. Eviction.} Standard KV cache eviction algorithms (H2O~\cite{h2023}, SnapKV~\cite{snapkv24}) drop cache entries to enforce a strict memory budget $C_{\text{max}}$. While effective at reducing peak VRAM, hard eviction suffers from irreversible forgetting. When an agent enters a subsequent reasoning turn requiring reference to an earlier code snippet or tool schema, the dropped context forces re-computation or induces hallucinations.

\noindent\textbf{Quantized KV Caches.} Quantization schemes such as KIVI~\cite{kivi24} compress key and value matrices to lower bit-widths. While lossless in gross retrieval tasks, ultra-low-bit representations (1-bit or 2-bit) experience accuracy drops in dense reasoning tasks unless critical outlier channels and attention sinks are preserved in higher precision.

\noindent\textbf{Draft-Based Approximate Inference.} SpecKV~\cite{speckv26} introduces the use of a smaller model from the same architectural family to generate a sequence of draft tokens. By computing cross-attention from draft queries back to prompt keys, SpecKV identifies influential tokens before performing eviction. However, SpecKV is limited to a single deterministic draft trace ($N=1$) and relies entirely on hard token dropping.

\noindent\textbf{Speculative Agent and Tree Execution.} Recent systems apply speculation across broader agent components. PASTE~\cite{sui2026parallelizingtoolexecutionllm} and Speculative Actions~\cite{specaction26} exploit predictable control and data flow patterns to speculatively execute external tool invocations and API calls, hiding round-trip environment latency. In tree-structured reasoning, SPEX~\cite{tot2026} breaks reward synchronization barriers by speculatively expanding promising branches in parallel. Our work is complementary: while these approaches speculate over external actions or reasoning frontiers, we apply speculative multi-branch consensus internally to manage the physical KV cache memory tier.

\section{Proposed Methodology}

\subsection{Triggered Multi Rollout Lookahead}
Running continuous speculative rollouts introduces compute overhead. We treat speculative compression as an elastic resource allocation mechanism, invoked only under high GPU memory pressure.

When triggered, the serving system forks the current sequence context to a lightweight draft model (e.g., Llama-3.2-1B for a Llama-3.1-8B/70B target). The draft model executes $N$ parallel stochastic rollouts (where $N \in [4, 8]$) for a lookahead horizon of $L = 10$ tokens. Because the draft model has a substantially smaller hidden dimension and runs with batched decoding, generating short speculative tokens introduces minimal latency.

\subsection{Consensus Attention Scoring}
For each rollout branch $j \in \{1, \dots, N\}$, we capture the cross-attention activations $A^{(j)} \in \mathbb{R}^{L \times T_{\text{ctx}}}$ between the draft query tokens and the existing context keys. We define the consensus importance score $S_i$ for context token $i$ as:
\begin{equation}
S_i = \frac{1}{N} \sum_{j=1}^N \max_{t \in \{1, \dots, L\}} A^{(j)}_{t, i}
\end{equation}
By computing the maximum attention across the horizon and averaging across the branches, $S_i$ highlights tokens that consistently matter across multiple probable reasoning paths, filtering out branch-specific noise.

\subsection{Dynamic Mixed-Precision Quantization}
The scheduler then partitions the KV cache as per the following:
\begin{itemize}
    \item \textbf{High-Precision Subset ($S_{\text{high}}$):} The top-$k\%$ tokens with the highest consensus scores $S_i$ (along with initial attention sinks and recent window tokens) are retained in FP16/BF16.
    \item \textbf{Compressed Subset ($S_{\text{low}}$):} The remaining tokens are dynamically quantized to 4-bit or 2-bit representations using channel-wise scaling factors.
\end{itemize}
This ensures that memory is aggressively freed while retaining structural information across the entire sequence.
\bibliographystyle{plain}
\bibliography{references}

\end{document}
